% Avsnitt 12.3, ur lectures/L12/README.md.

\section{Sammanfattning}
\label{c12:sec:review}

\needspace{6\baselineskip}
\subsection*{Det här ska du kunna nu}
\begin{itemize}
  \item Implementera en breddgenerisk dubbelvippsynkroniserare, och förklara varför den inte behöver
    någon reset och vad en breddning \emph{inte} löser.
  \item Implementera en räknarbaserad bittimer som producerar rätt tajmade \code{sample}- och
    \code{bit_done}-pulser, med en \code{resync}-ingång.
  \item Förklara vad \code{resync} gör och varför \code{sample} kommer före \code{bit_done}.
  \item Läsa ett spår på tickvis nivå över en bitperiod och för varje tick säga vad varje utgång bör
    vara.
  \item Köra en utdelad testbänk hela vägen och spåra ett fel tillbaka till den tick i designen som
    orsakade det.
\end{itemize}

\needspace{6\baselineskip}
\subsection*{Frågor att testa dig själv med}
\begin{enumerate}
  \item Varför behöver en synkroniserare två vippor snarare än en, och varför behöver den ingen
    reset?
  \item En breddad \code{meta_prev} tar bort metastabilitet men inte skevhet mellan bitarna. Vad är
    skillnaden, och när slutar en synkroniserare vara rätt verktyg?
  \item Vad gör \code{bit_timer.resync}, och när används den?
  \item Varför kommer \code{sample} före \code{bit_done} snarare än tvärtom?
  \item \code{bit_timer} håller sin räknare medan \code{enable = '0'}. Vilket tillstånd i
    \code{can_controller} är det enda som någonsin skulle vilja ha det, och varför visar det sig inte
    vilja det heller?
  \item När en utdelad testbänk misslyckas, vad skriver GHDL ut, och hur hittar du vilken kontroll
    som fallerade?
\end{enumerate}

\needspace{6\baselineskip}
\subsection*{Riktvärde}
\code{meta_prev} och \code{bit_timer} bör vara klara och mergade inom ungefär en vecka. Båda är
korta, och de är de enda modulerna som ingen annan modul kan ersätta. Det här är också de första
riktiga Pull Requesterna: använd granskningen, och kontrollera portordning, att testbänken passerar,
och att inga konstanter är hårdkodade förbi \code{can_def}.

\needspace{6\baselineskip}
\subsection*{Blicken framåt}
\Chapref{c13:ch} bygger \code{crc15}, en bitseriell checksummemotor som både genererar och
kontrollerar en CRC med samma logik och utan lägesomkoppling.
